\chapter{Subgap spectroscopy and Andreev qubits}\label{ch:experiments}

\section{From an eigenstate to a measured signal}\label{sec:measurement-picture}

A calculated energy difference determines where a resonance \emph{can} occur. An
experiment also specifies how the system is prepared, which operator couples to
the probe, and which response is detected. The same junction can therefore give
different tunnelling, microwave and supercurrent spectra. A useful procedure
for interpreting a measurement is
\begin{enumerate}
\item identify the initial state's energy, parity, spin and population;
\item identify the probe operator and its symmetry;
\item find final states with nonzero matrix elements;
\item compare their energy differences with the probe energy;
\item include temperature, broadening and the detector response.
\end{enumerate}
The last step often requires physics beyond the isolated-junction Hamiltonian.
An eigenstate calculation provides the energies and matrix elements that enter
the response calculation; a measured line shape also depends on the probe and environment.

\begin{figure}[htb]
\centering\manualfigure{primer_probes}
\caption{Two measurements of the same subgap system. A weak normal tunnel probe
exchanges one electron with the junction and measures a conductance response.
A microwave line modulates a gate or phase and is detected through absorption
or a resonator response. The associated parity-changing and parity-preserving
transitions have frequencies and visibility set by their respective energies, matrix elements and populations.}
\label{fig:primer-probes}
\end{figure}

\subsection{Single-electron tunnelling}

For an initial ground state $\ket g$, define excitation energies
$\omega_m=E_m-E_g\ge0$. The zero-temperature local spectral function is
\begin{equation}\label{eq:primer-spectral}
 A_\sigma(E)=\sum_m\left[
 |\matel m{d_\sigma\dg}g|^2\delta(E-\omega_m)
 +|\matel m{d_\sigma}g|^2\delta(E+\omega_m)\right].
\end{equation}
The first term adds an electron and the second removes one. The negative-bias
peak represents electron removal with the positive excitation cost $\omega_m$.
Both operators change the parity of the probed dot--superconductor subsystem.
The total parity including the probe remains conserved.

For a sufficiently weak, featureless normal probe in equilibrium spectroscopy,
$dI/dV$ samples a thermally broadened version of this spectral function. We use
$eV$ as the probe electron energy relative to the junction; experimental voltage
conventions must be translated before comparing signs. Peaks at opposite bias
can have different intensities because the electron and hole matrix elements
generally differ away from particle--hole symmetry.
Spin selection rules and occupation of the initial doublet component can suppress
one branch altogether.

The probe and preparation determine how this spectral function appears experimentally. A superconducting tip supplies
its own energy-dependent DOS and gap, shifting and reshaping thresholds. A normal
lead with appreciable coupling adds a gapless channel and can change the impurity
state it is meant to probe. Nonequilibrium occupation requires spectral contributions
from the populated states. Transport in an N--QD--S device can involve a complete
Andreev cycle transferring a pair even though its individual spectroscopic
steps exchange single electrons.

\subsection{Microwave absorption and dispersive sensing}

A weak modulation $\delta\lambda\cos(2\pi ft)$ couples through
$O=\partial H/\partial\lambda$. Gate modulation is associated with charge;
phase modulation with the current operator. These operators are parity even.
For a transition $i\to f$, the resonance condition is
\begin{equation}
 hf=E_f-E_i,
\end{equation}
and weak net absorption contains the factor
$(p_i-p_f)|\matel fOi|^2$, where $p_i$ and $p_f$ are populations. Observing a net
absorption line therefore requires both a nonzero matrix element and a population
contrast between the initial and final states.

A resonator can also detect the reactive response of a junction while being
detuned from its transitions. Virtual transitions shift the resonator's
frequency and change the phase of reflected or transmitted microwaves. This is
\emph{dispersive readout}. In a low-frequency adiabatic limit, the response is
related to the curvature of a state energy with phase, or its Josephson
inductance. At finite frequency it generally depends on matrix elements and
transition denominators as well~\cite{metzger2021circuitqed}. Interpret the frequency
shift through this full junction--resonator response.

\section{Superconducting quantum point contacts and atomic contacts}

The normal conductance of spin-degenerate channels is
\begin{equation}
 G_N=\frac{2e^2}{h}\sum_n\tau_n.
\end{equation}
It determines the sum of transmissions. Resolving the individual channels requires
the additional information in the nonlinear subgap current--voltage characteristic. Under voltage
bias, repeated Andreev reflections allow an excitation to gain energy in steps
of $eV$ until it reaches a continuum. For equal gaps, features appear near
$e|V|=2\Delta/n$, with positive integer $n$, and their amplitudes depend on
the channel transmissions. This \emph{multiple Andreev reflection} (MAR)
measurement is distinct from the equilibrium phase-biased spectrum of
\cref{eq:contact-abs}.

Scheer and collaborators used superconducting transport to determine the
transmission channels of atomic contacts and connect them to chemical
valence~\cite{scheer1998signature}. Thus a few-channel junction
has microscopic structure beyond its normal resistance. Once transmissions
are known, the short-contact formula predicts the corresponding phase-dependent
ABS branches under its assumptions.

Bretheau and collaborators detected the Andreev pair excitation of a
superconducting atomic contact using a voltage-biased Josephson junction as a
microwave spectrometer~\cite{bretheau2013exciting}. The spectrometer's Josephson
frequency satisfies $hf=2eV_J$, while the contact absorbs at $hf=2E_A$.
There are two different junctions and two different roles in this relation:
the voltage-biased source produces radiation, and the phase-biased contact
absorbs it. Identifying which junction supplies the bias keeps the factor of two
discussed in \cref{sec:abs-occupations} explicit.

The same concepts apply to gate-defined semiconductor constrictions.
Microwave sensing of a gate-defined superconducting QPC was demonstrated by
Chidambaram and collaborators~\cite{chidambaram2022microwave}. Semiconductor
junctions also permit gate-dependent channel occupation, appreciable spin--orbit
coupling and longer dwell times. These features make them tunable, but require
checking whether the simplest short, spin-degenerate channel formula is adequate.
Nanowire microwave spectroscopy provides complementary
examples~\cite{vanwoerkom2017microwave,tosi2019spinorbit}.

\section{Reading a quantum-dot tunnelling spectrum}

Deacon and collaborators studied a dot strongly coupled to a superconductor
and weakly coupled to a normal probe~\cite{deacon2010tunneling}. Gate-dependent
subgap resonances revealed the competition between singlet and doublet ground
states. Lee and collaborators resolved the field dependence and spin splitting
of these excitations~\cite{lee2014spinresolved}. These measurements connect the
many-body energies of \cref{sec:subgap} to experimental lines.

\begin{figure}[htb]
\centering\manualfigure{primer_spectroscopy}
\caption{Analytic interpretation exercise for the large-gap dot, with parameters
specified in \cref{fig:primer-atomic}. (a) Signed energy difference $E_D-E_S$.
(b) Positive and negative energies of the lowest single-electron transition,
$\pm|E_D-E_S|$. The curves specify resonance positions; a conductance calculation
would additionally require line widths, populations, transition weights and higher transitions.
The doublet is the ground state between the crossings.}
\label{fig:primer-spectroscopy}
\end{figure}

Start with a singlet ground state and an excited spin-degenerate doublet. A
small field splits the final doublet, so two transition energies can emerge.
On the doublet-ground-state side, sufficiently low temperature can populate
only the lower Zeeman component. Tunnelling to a singlet then gives a different
number of branches. Thermal population of the upper component can restore additional
lines. The number of visible lines thus reflects both the spectrum and the
preparation of the initial state.

A zero-bias crossing can occur when an ordinary singlet and doublet exchange
order. Distinguishing this ordinary parity transition from a Majorana mode
requires additional evidence. To interpret an apparent avoided crossing as
coherent hybridization, resolve the relevant line widths, thermal populations
and contributing transitions, and check which levels the Hamiltonian couples.

For a useful comparison with a paper, identify the lead gaps, gate-to-energy
conversion, total versus per-lead hybridization, field direction, probe coupling
and temperature. Match the observable as well as the Hamiltonian. Comparing a
calculated energy branch with a bright experimental ridge is meaningful only
after checking which transition produces that ridge.

\section{The even-parity Andreev-level qubit}\label{sec:pair-qubit}

The two even states of a short channel form the Andreev-level, or Andreev-pair,
qubit: $\ket g$ and $\ket e$ of \cref{eq:andreev-occupations}. Their transition
frequency is $2E_A/h$ in that model. A gate or phase modulation can drive it
when the corresponding off-diagonal matrix element is nonzero. Quantitative
control requires the wavefunctions or a suitable effective Hamiltonian together
with the phase dependence of $E_A$.

Janvier and collaborators demonstrated coherent manipulation of Andreev states
in superconducting atomic contacts, with resonator-based readout, Rabi
oscillations and Ramsey measurements~\cite{janvier2015coherent}. Rabi
oscillations show controlled population transfer; Ramsey fringes measure the
phase accumulated by a superposition. These go beyond observing an absorption
resonance. The odd sector is outside the chosen even qubit space: trapping one
quasiparticle can interrupt the operation and change the readout response.

\begin{figure}[htb]
\centering\manualfigure{primer_qubits}
\caption{Two distinct Andreev-qubit encodings. The pair qubit uses the even
ground and excited states; the spin qubit uses two states within an odd
manifold. The vertical scales are schematic and independent. Parity switching
leaves the selected computational space in either case. Spin labels denote
pseudospins when spin--orbit coupling mixes physical spin.}
\label{fig:primer-qubits}
\end{figure}

For an interacting dot, the corresponding even transition is $S\to S^*$.
Its frequency is $(E_{S^*}-E_S)/h$ and must be calculated independently of the lowest
parity-changing excitation. The large-gap pseudospin Hamiltonian makes the
distinct energy differences explicit. The driven two-level example in \cref{ch:ex-tls}
explains the associated Rabi physics; \cref{ch:ex-phase,ch:ex-gate} address
phase and gate driving in the finite-gap junction.

\section{The odd-parity Andreev spin qubit}\label{sec:spin-qubit}

An Andreev spin qubit (ASQ) uses two states within an odd-parity manifold. In a
spin-conserving model these are $D_\up$ and $D_\dn$. With spin--orbit coupling
they are more generally pseudospin states with mixed spin and orbital character.
The qubit splitting is the energy difference between these two odd-parity states.
Quasiparticle addition and the pair excitation $2E_A$ involve different transitions.

An odd occupation can arise in two ways. A quasiparticle can be trapped
in an otherwise non-interacting weak link; the duration of this occupation is then
set by the parity lifetime. Alternatively, Coulomb interaction can make the odd doublet the
equilibrium ground manifold over a range of gates and phases. The second route
connects the superconducting Anderson model directly to a qubit architecture.
Both require understanding population leakage and the gap to other states.

Hays and collaborators demonstrated continuous monitoring of a trapped spin
through its spin-dependent microwave response~\cite{hays2020continuous}. Their
subsequent coherent-control experiment used two-tone Raman transitions through
auxiliary Andreev states~\cite{hays2021coherent}. A Raman process couples the
two qubit states indirectly through driven transitions to higher levels. The
energies and matrix elements of those auxiliary levels are essential to the
control mechanism, even though the qubit itself is encoded in two states.

In quantum-dot devices, Bargerbos and collaborators measured spin-split
Andreev levels and spin-conditioned circuit spectra~\cite{bargerbos2023spectroscopy}.
Pita-Vidal and collaborators demonstrated direct electrical spin control and
strong coupling to a transmon~\cite{pitavidal2023direct}. Spectroscopic avoided
crossings establish hybridization between the spin and circuit excitations;
determining two-qubit gate fidelity additionally requires coherent-control measurements.

\subsection{Why can an electric field rotate a spin?}

A scalar gate modulation preserves an exactly conserved spin. Spin--orbit
coupling mixes spin and orbital motion so that an electric perturbation can
acquire a transverse matrix element in the qubit subspace. Assess the static
spin splitting and the transverse electrical coupling separately. A phase-biased
structure with suitable spin--orbit and orbital structure can produce spin splitting
at zero Zeeman field away from time-reversal-invariant operating points.
Spin--orbit coupling alone preserves Kramers degeneracy at those points, while
a spin-rotation-invariant junction retains spin degeneracy under phase bias.

This is a reason to distinguish a minimal single-orbital, spin-conserving model
from a realistic all-electric ASQ Hamiltonian. Spin-dependent tunnelling,
multiple orbitals or an additional conduction path may be essential. The
dot-local transverse magnetic drive of \cref{ch:ex-esr} isolates spin
renormalization and resonance physics; electrical ASQ control requires the
appropriate microscopic operator and symmetry breaking.

\section{Embedding the junction in a quantum circuit}\label{sec:circuit-picture}

So far $\phi$ has been an externally prescribed parameter. In a circuit,
capacitance introduces charging energy and makes phase a quantum variable
conjugate to transferred Cooper-pair number, $[\hat\phi,\hat n]=\ii$. For an
inductively shunted junction restricted to an isolated internal branch $p$, an
illustrative Hamiltonian is
\begin{equation}\label{eq:primer-circuit}
 H_{\mathrm{circ}}=4E_C(\hat n-n_g)^2+E_p(\hat\phi)
                  +\frac{E_L}{2}(\hat\phi-\phi_{\mathrm{ext}})^2,
 \qquad E_C=\frac{e^2}{2C}.
\end{equation}
The dimensionless $n_g$ is an offset charge in Cooper-pair units, and
$E_L=(\hbar/2e)^2/L$ for inductance $L$. A reference Josephson junction provides
a cosine contribution instead of this ideal inductive term. The circuit layout
determines which potential is appropriate.

\begin{figure}[htb]
\centering\manualfigure{primer_circuit}
\caption{Conceptual microwave experiment. A weak link and a reference junction
form a flux-controlled loop; a shunt capacitance and a coupled resonator supply
collective circuit degrees of freedom and readout. The sequence separates
preparation, coherent driving and measurement. Resonator output depends on the
occupied junction state and the coupled circuit response. The schematic illustrates
the general architecture and readout sequence.}
\label{fig:primer-circuit}
\end{figure}

Near a stable minimum of the total potential $V(\phi)$, small phase oscillations
have the approximate spacing $\hbar\omega\simeq\sqrt{8E_CV''(\phi_{\min})}$.
This spacing follows from the circuit capacitance and potential curvature.
The branch projection in
\cref{eq:primer-circuit} is adiabatic: near internal-state resonances one must
retain a matrix-valued junction Hamiltonian and its couplings. Charge offsets
and boundary conditions must also be consistent with the chosen parity.
Reviews in Refs.~\cite{aguado2020perspective,glazman2021bogoliubov} explain this
connection between microscopic states and collective circuit modes.

\section{Lifetimes, poisoning and the limits of a level diagram}

Three distinct processes are relevant here. \emph{Relaxation} transfers population
between levels within the chosen system. \emph{Dephasing} destroys relative
phase coherence of a superposition. \emph{Poisoning} changes the parity of the
bound subsystem by exchanging a quasiparticle with its environment. A spin flip
need not change parity; a parity switch need not be a spin flip. Escape of a
driven quasiparticle into the continuum can empty a bound state while preserving
the total parity of bound state plus reservoirs.

For a non-interacting channel, promoting an occupied bound quasiparticle to the
continuum has threshold $\Delta-E_A$, while creating one bound and one continuum
quasiparticle costs $\Delta+E_A$. The even pair transition costs $2E_A$.
Identify the relevant final states and excitation threshold for each microwave
process. Interacting thresholds use the appropriate
many-body final energies; \cref{ch:ex-phase} gives a concrete example.

Measured $T_1$, $T_2$ and line widths depend on electromagnetic noise,
quasiparticle distributions, drive strength and couplings to external modes.
Calculating these rates requires a model of the relevant dynamical processes.
Numerical Lorentzian broadening specifies the plotted line shape; interpreting
its width as a lifetime requires that physical model. The package's states and
operator matrix elements supply ingredients for calculating those responses.

\begin{physics}[Read the experiment]
A resonator trace switches between two frequencies. Which state changes could
produce this signal? It may resolve populations switching between spin states or
between parities. Identify the state-dependent response and the switching
process. To test coherent spin control, use controlled oscillations under a pulse and phase-sensitive
Ramsey measurements. Similarly, an absorption line
locates a transition, while its disappearance can reflect population changes,
a selection rule or escape from the bound manifold.
\end{physics}

\paragraph{A route through the literature.}
For devices and transport, start with De Franceschi and
collaborators~\cite{defranceschi2010hybrid}; for dot-junction theory, use
Mart\'in-Rodero and Levy Yeyati~\cite{martinrodero2011josephson} and
Meden~\cite{meden2019andersonjosephson}. Read the atomic-contact spectroscopy
and coherent-control papers together~\cite{bretheau2013exciting,janvier2015coherent}
to see the progression from an excitation to a qubit. Then compare trapped-spin
control~\cite{hays2021coherent} with the Coulomb-stabilized dot
implementation~\cite{pitavidal2023direct}. This order keeps the encoding,
measurement and microscopic model distinct while showing how they fit together.
